\documentclass[10pt,a4paper]{amsart}
\usepackage[margin=19mm]{geometry}
\usepackage{amsmath,amssymb,mathtools}
\usepackage[scr=rsfs,cal=euler]{mathalfa}
\usepackage{xfrac}
\usepackage[unicode,hidelinks,bookmarks=false]{hyperref}
\newcommand{\ie}{i.e.\ }
\newcommand{\cf}{cf.\ }
\newcommand{\resp}{respectively\ }
\newcommand{\Subsection}{\subsection}
\providecommand{\nobreakdash}{\nobreak\hbox{-}}
\setlength{\emergencystretch}{2em}
\multlinegap0pt
\usepackage{amssymb}
\usepackage[shortlabels]{enumitem}
\usepackage{bm}
\usepackage{xcolor}
\usepackage{mathtools}
\usepackage{dsfont}
\newtheorem{theorem}{Theorem}
\newtheorem{corollary}{Corollary}
\usepackage{cleveref}
\crefname{theorem}{Theorem}{Theorems}
\crefname{corollary}{Corollary}{Corollarys}
\def\NN{\mathbb N}
\renewcommand{\geq}{\geqslant}
\renewcommand{\leq}{\leqslant}
\title{Are trees really just butterflies in disguise?}
\author{Giovanne Santos, Maya Stein, Ella Williams}
\date{Source: arXiv:2609.09142v2 (CC BY 4.0)}
\begin{document}
\maketitle

\noindent\emph{Source and licence.} Are trees really just butterflies in disguise?. Version arXiv:2609.09142v2 (CC BY 4.0), distributed by arXiv under the Creative Commons Attribution 4.0 International licence, http://creativecommons.org/licenses/by/4.0/, as recorded in the arXiv licence metadata for this exact version and shown on the abstract page https://arxiv.org/abs/2609.09142v2. Full licence notice: \texttt{licenses/arxiv-2609.09142-CC-BY-4.0.txt}.\par\smallskip

\noindent\textit{Editorial scope.} Counted unit: the article's corollary corollary (source0.tex lines 720--725). The article's complete original proof of it is delivered with it (source0.tex lines 727--734, one \texttt{proof} environment of 8 lines). No mathematics is written, repaired or re-derived: the statement and the proof are the article's own lines, in the article's own notation, and no retained line is substituted or re-flowed.  Retained uncounted as the source-local context the proof needs: lines 188--190.  Nothing was omitted from any retained span, so the package carries no omission record.  No retained line contains a citation, so the wrapper carries no bibliography and no bibliography entry.  No retained line contains a figure environment, an image-inclusion command or drawing code, so no image is delivered and the package declares no asset dependency.  Editorial formatting only, in addition: an AMS \texttt{amsart} preamble that restates the article's own declarations and macros used by the retained lines, namely the environments \texttt{theorem}, \texttt{corollary} on separate counters, together with the standard packages those lines need.  The retained environments print in this document as Theorem 1, Corollary 1 in order of appearance, whereas the article prints the same items with the numbering of its own full document.  The complete native source of the article as distributed in the arXiv e-print for this exact version is bundled unchanged as \texttt{sources/209782/source0.tex}.
\begin{theorem}\label{thm:main_dense}
    For all $\alpha>0$ and $\Delta\in \NN$ there exists $k_0$ such that the following holds for all $n\geq k>k_0$ with $k\geq \alpha n$. Every digraph  $D$ on $n$ vertices with $e(D)\geq (1+\alpha)kn$ contains a copy of every antidirected tree $T$ on $k$ arcs with $\Delta(T)\leq \Delta$.
\end{theorem}

\begin{corollary}
    For all $\alpha>0$ and $\Delta,\ell \in \NN$ with $\ell\geq 2$ there exists $k_0$ such that for all $k>k_0$ the following holds. For every antidirected tree $T$ on $k$ arcs with $\Delta(T)\leq \Delta$, we have 
    \begin{equation*}
        \overleftrightarrow{r}(T,\ell)\leq (1+\alpha)\ell k \quad \text{ and } \quad \overrightarrow r(T,\ell)\le (2+\alpha)\ell k.
    \end{equation*}
\end{corollary}

\begin{proof}
    Let $\alpha>0$, $\Delta\in \NN$ and $\ell \geq 2$.
    We may assume that $\alpha\in (0,1)$ is sufficiently small so that $\alpha\leq 1/(4\ell)$ holds, as this only strengthens the statement. Let $k_0$ be the output of~\cref{thm:main_dense} with inputs $\alpha$ and~$\Delta$.  
    We first prove $\overleftrightarrow{r}(T,\ell)\leq (1+\alpha)\ell k$.
    Let $D$ be an $\ell$-edge-coloured complete digraph on $n = (1+\alpha)\ell k$ vertices, and note that $\alpha n \leq 2\alpha \ell k \leq k$. By averaging, there is a colour $c$ that appears on at least $n(n-1)/\ell = (1+\alpha)k(n-1)> (1+\alpha/2)kn$ arcs. Applying \cref{thm:main_dense} to the spanning subdigraph consisting of all arcs of colour $c$, we find a monochromatic copy of $T$ in $D$, as desired.

    The same averaging argument, together with the observation that every tournament contains exactly half of the arcs of the complete symmetric digraph, yields the proof of $\overrightarrow r(T,\ell)\le (2+\alpha)\ell k$.
\end{proof}

\end{document}
